% BEGIN REV09_NIVELES_2_5
\subsection{Compatibilidad de los lectores modulares de niveles dos y cinco}
\label{corpus9:iv:niveles}

Sobre la realización modular posterior a la construcción HMT,
fijemos \(\operatorname{Im}\tau>0\), \(q=e^{2\pi i\tau}\) y
\(q^{1/5}=e^{2\pi i\tau/5}\). El lector de Rogers--Ramanujan es
\[
 R(q)=\frac{q^{1/5}}{1+\dfrac{q}{1+\dfrac{q^2}{1+\cdots}}},
 \qquad x=R(q)^5.
\]
La identidad modular de esta realización, con la normalización
de \(j\) usada en \eqref{km:j-nivel2}, es
\begin{equation}
 j=-\frac{(x^4-228x^3+494x^2+228x+1)^3}
          {x(x^2+11x-1)^5}.
 \label{corpus9:iv:nivel5-identidad}
\end{equation}
La identidad constituye la premisa de realización modular del
siguiente transporte algebraico; las semillas, la orientación y
las representaciones de autoescala permanecen las de APP--TRIT--TPK.

\begin{proposition}[Cociente pentádico y selección de rama]
\label{corpus9:iv:nivel5-cociente}
Para \(x\ne0\), sea \(u=x-x^{-1}\). En el dominio \(u\ne-11\),
\begin{equation}
 j=-\frac{(u^2-228u+496)^3}{(u+11)^5}.
 \label{corpus9:iv:nivel5-u}
\end{equation}
La involución \(x\mapsto-1/x\) conserva \(u\). Sus preimágenes son
\[
 x=\frac{u\pm\sqrt{u^2+4}}2.
\]
El recubrimiento cuadrático se ramifica en \(u=\pm2i\); la
aplicación racional posterior \(u\mapsto j\) tiene grado seis.
\end{proposition}
\begin{proof}
Las factorizaciones exactas son
\[
 \begin{aligned}
 x^4-228x^3+494x^2+228x+1
   &=x^2(u^2-228u+496),\\
 x^2+11x-1&=x(u+11).
 \end{aligned}
\]
Su sustitución en \eqref{corpus9:iv:nivel5-identidad} prueba
\eqref{corpus9:iv:nivel5-u}. La ecuación de la fibra es
\(x^2-ux-1=0\); su discriminante da los puntos de ramificación.
El numerador de \eqref{corpus9:iv:nivel5-u} tiene grado seis
y el denominador grado cinco. En \(u=-11\), el polinomio
del numerador vale \(3125\), de modo que ambos carecen de
factor común. Esto prueba el grado indicado.
\end{proof}

Para \(u\) real, las dos raíces tienen signos opuestos. En el
dominio complejo se conserva la rama elegida en la parametrización
por \(\tau\). La lectura \(j\) por sí sola conserva una fibra
genéricamente de seis valores de \(u\).

\begin{proposition}[Límite de autoescala y compatibilidad de niveles]
\label{corpus9:iv:niveles-ramas}
En el límite real \(q\to1^-\),
\[
 R(q)\longrightarrow\varphi_{\mathrm{HMT}}^{-1},\qquad
 x\longrightarrow\varphi_{\mathrm{HMT}}^{-5}
     =\frac{5\sqrt5-11}{2},\qquad u\longrightarrow-11.
\]
Cuando \(t=t_2(\tau)\) y \(u\) se leen sobre el mismo parámetro
modular, satisfacen
\begin{equation}
 (t+256)^3(u+11)^5
   +t^2(u^2-228u+496)^3=0.
 \label{corpus9:iv:niveles-enlace}
\end{equation}
Poniendo \(\epsilon=u+11\), las ramas que cumplen los límites
indicados tienen las leyes
\begin{align}
 \epsilon\to0,\ t\to0:
 &\quad \frac{t^2}{(-\epsilon)^5}\longrightarrow
          \frac{2^{24}}{5^{15}},\\
 \epsilon\to0,\ t\to\infty:
 &\quad t\epsilon^5\longrightarrow-5^{15}.
 \label{corpus9:iv:niveles-asintotica}
\end{align}
\end{proposition}
\begin{proof}
Para \(0<q<1\), los convergentes positivos alternos de la
fracción continua encierran \(R(q)\). Cada convergente fijo
tiende, cuando \(q\to1^-\), al convergente correspondiente de
la fracción continua con todos los numeradores iguales a uno.
Las dos subsucesiones de esos convergentes tienen límite
\(\varphi_{\mathrm{HMT}}^{-1}\), en su realización cuadrática
ya establecida. El encaje prueba el primer límite. Elevar a
la quinta potencia y usar \(\varphi^2=\varphi+1\) da los
otros dos.

Igualar \eqref{km:j-nivel2} y \eqref{corpus9:iv:nivel5-u}
prueba \eqref{corpus9:iv:niveles-enlace}. La sustitución
\(u=\epsilon-11\) da
\[
 \frac{(t+256)^3}{t^2}
   =-\frac{(3125-250\epsilon+\epsilon^2)^3}{\epsilon^5}.
\]
En la rama \(t\to0\), la identidad se reordena como
\[
 \frac{t^2}{(-\epsilon)^5}
   =\frac{(t+256)^3}{(3125-250\epsilon+\epsilon^2)^3}.
\]
El límite es \(256^3/3125^3=2^{24}/5^{15}\).
En la rama \(t\to\infty\), multiplicar por \(\epsilon^5\)
y escribir el miembro izquierdo como
\(t\epsilon^5(1+256/t)^3\) da el segundo límite.
\end{proof}

La relación conserva las ramas de una correspondencia sobre el
mismo \(j\). La elección de \(\tau\), o del levantamiento
correspondiente, fija la pareja que se transporta. El polo
pentádico y la inversión de Fricke mantienen sus dominios y sus
operaciones respectivas.
% END REV09_NIVELES_2_5
